\section{Discussion and Limitations}
\label{sec:discussion}

\stitle{What the results do and do not show.}
Memory is cheap where the full history is not, and the two tiers answer different
questions: the record holds detail and facts said once, the notes hold what is
current. The case for the notes, however, rests on extraction, and our rule-based
extractor finds only about half the facts in the benchmark. A language model behind
the same signature is the largest open gain, and the benchmark's oracle rows show
what it could deliver. We have not built or measured one.

\stitle{Retrieval.}
The index matches words, not meaning. Stemming and stopword removal fixed the most
visible failures, but a paraphrase that shares no words with the query is not found.
\sys accepts any embedding function, and testing it on a workload designed to defeat
lexical search is future work.

\stitle{Scale and realism.}
Both benchmarks are synthetic, at small scale, and run in seconds. We have not
evaluated on recorded conversations or with a model as the reader, and we make no
claim about accuracy on public long-conversation benchmarks. The harness takes a
workload and a retrieval function, so both are straightforward to add.

\stitle{Design costs.}
A single log with a single writer per store keeps commits ordered and guards
meaningful, at the price of a server when agents run in many processes. Cascading
invalidation is the right default, but a frequently updated, heavily cited note
stales a lot; a flag turns it off for minor edits. Conflicts resolve by supersession,
not by judgement: the newest statement wins until someone retracts the session it came
from.

\section{Conclusion}
\label{sec:conclusion}

Agent memory is a database problem, and a log with provenance gives it transactions,
invalidation, and rollback without special cases. On that base, two tiers of memory, an
index, tools, a declared structure, and an offline consolidation pass fit together
because each is a transaction with a record of where it came from. Measured honestly
on a synthetic benchmark, the record alone is a strong baseline, and the notes are as
good as the extractor behind them.
